%% Solutions for Chapter 8: Regular and Context-Free Grammars

\chapter{Regular and Context-Free Grammars --- Solutions}

\begin{enumerate}[{8.}1. ]

%%------------------------------------------------------------
\item \textbf{Can strings like $\pmb{ab}BA\pmb{d}B\pmb{c}$ be derived from the start symbol in a right-linear grammar?}

No. In a right-linear grammar, every production has the form $A \to w$ or $A \to wB$ where $w \in \Sigma^*$ and $B$ is a nonterminal. Thus any sentential form derived from the start symbol has the structure $w$ or $wB$ where all terminal symbols precede the single nonterminal (if any). A string like $\pmb{ab}BA\pmb{d}B\pmb{c}$ contains nonterminals interspersed with and following terminals, which cannot occur in any right-linear grammar derivation.

%%------------------------------------------------------------
\item \textbf{Prove $P(n)$: for all $x \in \Sigma^n$, if $t_0 \DRightarrow xt_j$ then $\DBAR_A(t_0, x) = t_j$.}

\textbf{Base case} ($n=0$): The only string of length 0 is $\lambda$. The derivation $t_0 \DRightarrow \lambda t_0$ uses zero steps, and $\DBAR_A(t_0, \lambda) = t_0$ by definition, so $P(0)$ holds.

\textbf{Inductive step}: Assume $P(n)$ holds for all strings of length $n$. Let $x = ya$ where $y \in \Sigma^n$ and $a \in \Sigma$. If $t_0 \DRightarrow ya \cdot t_j$, then the derivation must pass through some intermediate state: $t_0 \DRightarrow y \cdot t_k \Rightarrow ya \cdot t_j$ for some $t_k$. By the inductive hypothesis applied to $y$, $\DBAR_A(t_0, y) = t_k$. The final step $t_k \Rightarrow a \cdot t_j$ corresponds to a production $t_k \to \pmb{a} t_j$ in the grammar (by construction of $G_A$), which corresponds to a transition $\delta(t_k, a) = t_j$ in $A$. Therefore $\DBAR_A(t_0, ya) = \delta(\DBAR_A(t_0, y), a) = \delta(t_k, a) = t_j$.

%%------------------------------------------------------------
\item \textbf{Regular expressions for the languages of:}

\begin{enumerate}
\item $G_4$: Tracing the grammar, we get
\[
\begin{aligned}
S &= \pmb{ab}A \cup \pmb{bb}B \cup \pmb{cc}V,\\
A &= \pmb{b}C \cup \pmb{c}X,\qquad
B = \pmb{ab},\qquad
C = \lambda \cup \pmb{c}S,\\
V &= \pmb{a}V \cup \pmb{c}X,\qquad
X = \pmb{b}V \cup \pmb{aa}X.
\end{aligned}
\]
From
\[
V=\pmb{a}V\cup\pmb{c}X,
\]
Arden's lemma (Theorem 6.1) gives
\[
V=\pmb{a}^*\pmb{c}X.
\]
Substituting into the equation for $X$ gives
\[
X=\pmb{b}\pmb{a}^*\pmb{c}X\cup\pmb{aa}X
  =(\pmb{b}\pmb{a}^*\pmb{c}\cup\pmb{aa})X.
\]
There is no constant term, so the only solution is
\[
X=\emptyset,
\]
and therefore also $V=\emptyset$.

Thus only the $S,A,B,C$ portion contributes:
\[
S=\pmb{ab}A\cup\pmb{bbab},\qquad
A=\pmb{b}C\cup\pmb{c}\emptyset=\pmb{b}C,\qquad
C=\lambda\cup\pmb{c}S.
\]
Hence
\[
A=\pmb{b}\cup\pmb{bc}S,
\]
so
\[
S=\pmb{abb}\cup\pmb{abbc}S\cup\pmb{bbab}.
\]
By Arden's lemma (Theorem 6.1),
\[
L(G_4)=(\pmb{abbc})^*(\pmb{abb}\cup\pmb{bbab}).
\]

\item $G_5$: The language equations for the three nonterminals are
\[
\begin{aligned}
S_0 &= \lambda \cup \pmb{0}S_2 \cup \pmb{1}S_1,\\
S_1 &= \pmb{0}S_1 \cup \pmb{1}S_2,\\
S_2 &= \pmb{0}S_2 \cup \pmb{1}S_0.
\end{aligned}
\]
The equation for $S_2$ has the form required by Theorem 6.1, so
\[
S_2=\pmb{0}^*\pmb{1}S_0.
\]
Substituting this into the equation for $S_1$ gives
\[
S_1=\pmb{0}S_1\cup\pmb{1}\pmb{0}^*\pmb{1}S_0,
\]
and another application of Theorem 6.1 gives
\[
S_1=\pmb{0}^*\pmb{1}\pmb{0}^*\pmb{1}S_0.
\]
Now substitute both expressions into the equation for $S_0$:
\[
\begin{aligned}
S_0
&=\lambda\cup\pmb{0}\pmb{0}^*\pmb{1}S_0
  \cup \pmb{1}\pmb{0}^*\pmb{1}\pmb{0}^*\pmb{1}S_0\\
&=\lambda\cup(\pmb{0}^+\pmb{1}\cup
  \pmb{1}\pmb{0}^*\pmb{1}\pmb{0}^*\pmb{1})S_0.
\end{aligned}
\]
Thus
\[
L(G_5)=(\pmb{0}^+\pmb{1}\cup
\pmb{1}\pmb{0}^*\pmb{1}\pmb{0}^*\pmb{1})^*.
\]

\item $G_6$: From the productions $Z \to \pmb{a}Z | \pmb{b}T$, $T \to \pmb{a}Z$. Language equations: $Z = \pmb{a}Z \cup \pmb{b}T$ and $T = \pmb{a}Z$. Substituting: $Z = \pmb{a}Z \cup \pmb{ba}Z = (\pmb{a}\cup\pmb{ba})Z$. Since there is no terminal production (no production $Z \to w$ or $T \to w$ for $w \in \Sigma^*$), $L(G_6) = \emptyset$.

Every derivation cycles: $Z \Rightarrow \pmb{b}T \Rightarrow \pmb{ba}Z \Rightarrow \cdots$, and neither $Z$ nor $T$ has a terminal production, confirming $L(G_6) = \emptyset$.

\item $G_7$: $S \to \pmb{a}S | \pmb{ab}B | \pmb{c}C$, $B \to \pmb{ab}B | \lambda$, $C \to \pmb{c}C | \pmb{ca}$.
\begin{itemize}
\item From $B$: $B = \pmb{ab}B \cup \lambda$, so $L(B) = (\pmb{ab})^*$.
\item From $C$: $C = \pmb{c}C \cup \pmb{ca}$, so $L(C) = \pmb{c}^*\pmb{ca} = \pmb{c}^+\pmb{a}$.
\item From $S$: $S = \pmb{a}S \cup \pmb{ab}B \cup \pmb{c}C = \pmb{a}S \cup \pmb{ab}(\pmb{ab})^* \cup \pmb{c}^+\pmb{a}$.
  Thus
  \[
    S = \pmb{a}S \cup (\pmb{ab})^+ \cup \pmb{c}^+\pmb{a},
  \]
  so
  \[
    L(S) = \pmb{a}^*((\pmb{ab})^+ \cup \pmb{c}^+\pmb{a}).
  \]
\end{itemize}
The regular expression is $\pmb{a}^*(\pmb{ab})^+ \cup \pmb{a}^*\pmb{c}^+\pmb{a}$.
\end{enumerate}

%%------------------------------------------------------------
\item \textbf{Use Exercise 8.2 to formally prove Lemma 8.1.}

Lemma 8.1 states that for the grammar $G_A$ constructed from DFA $A$, $L(G_A) = L(A)$.

By the inductive result of Exercise 8.2, for all $x \in \Sigma^*$ and all states $t_i, t_j$: $t_i \DRightarrow x t_j$ in $G_A$ iff $\DBAR_A(t_i, x) = t_j$. Taking $t_i = t_0$ (start state):

\[
\begin{aligned}
x \in L(G_A)
&\Longleftrightarrow t_0 \DRightarrow x \text{ in } G_A\\
&\Longleftrightarrow (\exists t_j\in F)\ t_0 \DRightarrow x t_j \Rightarrow x\\
&\Longleftrightarrow (\exists t_j\in F)\ \DBAR_A(t_0,x)=t_j\\
&\Longleftrightarrow \DBAR_A(t_0,x)\in F\\
&\Longleftrightarrow x\in L(A).
\end{aligned}
\]
\qed

%%------------------------------------------------------------
\item \textbf{Draw the automata for the grammars in Exercise 8.3.}

Use the construction of Lemma~8.2: states are the suffix-strings that arise
from right-hand sides, the start state is $\langle S\rangle$, and the only
final state is $\langle\lambda\rangle$.

\begin{enumerate}
\item For $G_4$, the state set is
\[
\{\langle S\rangle,\langle A\rangle,\langle B\rangle,\langle C\rangle,\langle V\rangle,
\langle W\rangle,\langle X\rangle,\langle\pmb{ab}A\rangle,\langle\pmb{b}A\rangle,
\langle\pmb{bb}B\rangle,\langle\pmb{b}B\rangle,\langle\pmb{cc}V\rangle,
\langle\pmb{c}V\rangle,\langle\pmb{b}C\rangle,\langle\pmb{c}X\rangle,
\langle\pmb{ab}\rangle,\langle\pmb{b}\rangle,\langle\pmb{c}S\rangle,
\langle\pmb{a}V\rangle,\langle\pmb{aa}\rangle,\langle\pmb{a}W\rangle,
\langle\pmb{b}V\rangle,\langle\pmb{aa}X\rangle,\langle\pmb{a}X\rangle,
\langle\lambda\rangle\}.
\]
The nonempty $\lambda$-transitions are
\[
\delta(\langle S\rangle,\lambda)=\{\langle\pmb{ab}A\rangle,\langle\pmb{bb}B\rangle,
\langle\pmb{cc}V\rangle\},
\]
\[
\delta(\langle A\rangle,\lambda)=\{\langle\pmb{b}C\rangle,\langle\pmb{c}X\rangle\},
\quad
\delta(\langle B\rangle,\lambda)=\{\langle\pmb{ab}\rangle\},
\]
\[
\delta(\langle C\rangle,\lambda)=\{\langle\lambda\rangle,\langle\pmb{c}S\rangle\},
\quad
\delta(\langle V\rangle,\lambda)=\{\langle\pmb{a}V\rangle,\langle\pmb{c}X\rangle\},
\]
\[
\delta(\langle W\rangle,\lambda)=\{\langle\pmb{aa}\rangle,\langle\pmb{a}W\rangle\},
\quad
\delta(\langle X\rangle,\lambda)=\{\langle\pmb{b}V\rangle,\langle\pmb{aa}X\rangle\}.
\]
The nonempty letter-transitions are the suffix moves, for example
\[
\delta(\langle\pmb{ab}A\rangle,\pmb{a})=\{\langle\pmb{b}A\rangle\},\qquad
\delta(\langle\pmb{b}A\rangle,\pmb{b})=\{\langle A\rangle\},
\]
\[
\delta(\langle\pmb{ab}\rangle,\pmb{a})=\{\langle\pmb{b}\rangle\},\qquad
\delta(\langle\pmb{b}\rangle,\pmb{b})=\{\langle\lambda\rangle\},
\]
and the same suffix rule handles every other string-state: if the current state
is $\langle \mathbf{a}u\rangle$ and the next input letter is $\mathbf{a}$, the
machine moves to $\langle u\rangle$; if the input letter does not match the
first symbol, there is no transition. For example,
\[
\delta(\langle\pmb{bb}B\rangle,\pmb{b})=\{\langle\pmb{b}B\rangle\},\qquad
\delta(\langle\pmb{cc}V\rangle,\pmb{c})=\{\langle\pmb{c}V\rangle\},\qquad
\delta(\langle\pmb{aa}\rangle,\pmb{a})=\{\langle\pmb{a}\rangle\}.
\]

\item For $G_5$, the corresponding automaton is simply the three-state DFA
\[
\langle\{\pmb{0},\pmb{1}\},\{S_0,S_1,S_2\},S_0,\delta,\{S_0\}\rangle,
\]
with
\[
\delta(S_0,\pmb{0})=S_2,\quad \delta(S_0,\pmb{1})=S_1,\quad
\delta(S_1,\pmb{0})=S_1,\quad \delta(S_1,\pmb{1})=S_2,
\]
\[
\delta(S_2,\pmb{0})=S_2,\quad \delta(S_2,\pmb{1})=S_0.
\]

\item For $G_6$, the state set is
\[
\{\langle Z\rangle,\langle T\rangle,\langle\pmb{a}Z\rangle,\langle\pmb{b}T\rangle,
\langle\lambda\rangle\},
\]
with start state $\langle Z\rangle$, final state $\langle\lambda\rangle$, and
nonempty transitions
\[
\delta(\langle Z\rangle,\lambda)=\{\langle\pmb{a}Z\rangle,\langle\pmb{b}T\rangle\},
\qquad
\delta(\langle T\rangle,\lambda)=\{\langle\pmb{a}Z\rangle\},
\]
\[
\delta(\langle\pmb{a}Z\rangle,\pmb{a})=\{\langle Z\rangle\},\qquad
\delta(\langle\pmb{b}T\rangle,\pmb{b})=\{\langle T\rangle\}.
\]
There is no path to $\langle\lambda\rangle$, so the language is empty.

\item For $G_7$, the state set is
\[
\{\langle S\rangle,\langle B\rangle,\langle C\rangle,\langle\pmb{a}S\rangle,
\langle\pmb{ab}B\rangle,\langle\pmb{b}B\rangle,\langle\pmb{c}C\rangle,
\langle\pmb{ab}\rangle,\langle\pmb{b}\rangle,\langle\pmb{ca}\rangle,
\langle\pmb{a}\rangle,\langle\lambda\rangle\},
\]
with start state $\langle S\rangle$, final state $\langle\lambda\rangle$, and
nonempty $\lambda$-transitions
\[
\delta(\langle S\rangle,\lambda)=\{\langle\pmb{a}S\rangle,\langle\pmb{ab}B\rangle,
\langle\pmb{c}C\rangle\},
\]
\[
\delta(\langle B\rangle,\lambda)=\{\langle\pmb{ab}B\rangle,\langle\lambda\rangle\},
\qquad
\delta(\langle C\rangle,\lambda)=\{\langle\pmb{c}C\rangle,\langle\pmb{ca}\rangle\},
\]
and nonempty letter-transitions
\[
\delta(\langle\pmb{a}S\rangle,\pmb{a})=\{\langle S\rangle\},\quad
\delta(\langle\pmb{ab}B\rangle,\pmb{a})=\{\langle\pmb{b}B\rangle\},\quad
\delta(\langle\pmb{b}B\rangle,\pmb{b})=\{\langle B\rangle\},
\]
\[
\delta(\langle\pmb{c}C\rangle,\pmb{c})=\{\langle C\rangle\},\quad
\delta(\langle\pmb{ab}\rangle,\pmb{a})=\{\langle\pmb{b}\rangle\},\quad
\delta(\langle\pmb{b}\rangle,\pmb{b})=\{\langle\lambda\rangle\},
\]
\[
\delta(\langle\pmb{ca}\rangle,\pmb{c})=\{\langle\pmb{a}\rangle\},\qquad
\delta(\langle\pmb{a}\rangle,\pmb{a})=\{\langle\lambda\rangle\}.
\]
\end{enumerate}

%%------------------------------------------------------------
\item \textbf{Right-linear grammars for the given languages.}

\begin{enumerate}
\item All words in $\{\pmb{a},\pmb{b},\pmb{c}\}^*$ \emph{not} containing two consecutive $\pmb{b}$s:
\[
G = \langle\{S,A\},\{\pmb{a},\pmb{b},\pmb{c}\},S,\{S\to\pmb{a}S|\pmb{b}A|\pmb{c}S|\lambda,
A\to\pmb{a}S|\pmb{c}S|\lambda\}\rangle
\]
State $S$ = ``last symbol was not $\pmb{b}$'' (or start); state $A$ = ``last symbol was $\pmb{b}$''. From $A$, reading another $\pmb{b}$ is disallowed.

\item All words \emph{containing} two consecutive $\pmb{b}$s:
\[
G = \langle\{S,A,B\},\{\pmb{a},\pmb{b},\pmb{c}\},S,\{S\to\pmb{a}S|\pmb{c}S|\pmb{b}A,
A\to\pmb{b}B|\pmb{a}S|\pmb{c}S,B\to\pmb{a}B|\pmb{b}B|\pmb{c}B|\lambda\}\rangle
\]
State $S$ means ``no consecutive $\pmb{b}$s have been seen and the last symbol
was not $\pmb{b}$''; state $A$ means ``no consecutive $\pmb{b}$s have been seen,
but the last symbol was $\pmb{b}$''; state $B$ means ``two consecutive
$\pmb{b}$s have occurred.''  Once $B$ is reached, any remaining string may be
generated.

\item All words with $|w|_{\pmb{a}} = |w|_{\pmb{b}}$: This language is not regular (it includes $\{\pmb{a}^n\pmb{b}^n\}$), so no right-linear grammar exists.

\item All words with an even number of $\pmb{a}$s:
\[
G = \langle\{S,A\},\{\pmb{a},\pmb{b},\pmb{c}\},S,\{S\to\pmb{b}S|\pmb{c}S|\pmb{a}A|\lambda,
A\to\pmb{b}A|\pmb{c}A|\pmb{a}S\}\rangle
\]
State $S$ records that an even number of $\pmb{a}$s has been generated so far,
and state $A$ records that an odd number has been generated.  The
$\pmb{b}$- and $\pmb{c}$-productions preserve parity, while each $\pmb{a}$
switches parity.

\item All words \emph{not} ending in $\pmb{b}$:
\[
G = \langle\{S,A\},\{\pmb{a},\pmb{b},\pmb{c}\},S,\{S\to\pmb{a}S|\pmb{b}A|\pmb{c}S|\pmb{a}|\pmb{c}|\lambda,
A\to\pmb{a}S|\pmb{b}A|\pmb{c}S|\pmb{a}|\pmb{c}\}\rangle
\]
(Equivalently, $G = \langle\{S,A\},\ldots,\{S\to(\pmb{a}|\pmb{b}|\pmb{c})^*(\pmb{a}|\pmb{c})|\lambda\}\rangle$.)
State $S$ records that the string generated so far is empty or currently ends
in a symbol other than $\pmb{b}$; state $A$ records that it currently ends in
$\pmb{b}$.  Only $S$ may stop by $\lambda$, and both states may stop after
adding one final $\pmb{a}$ or $\pmb{c}$.

\item All words containing no $\pmb{c}$s:
\[
G = \langle\{S\},\{\pmb{a},\pmb{b},\pmb{c}\},S,\{S\to\pmb{a}S|\pmb{b}S|\lambda\}\rangle
\]
\end{enumerate}

%%------------------------------------------------------------
\item \textbf{Left-linear grammars for the languages in Exercise 8.6.}

A left-linear grammar has productions of the form $A \to w$ or $A \to Bw$.

\begin{enumerate}
\item No two consecutive $\pmb{b}$s:
\[
G = \langle\{S,A\},\{\pmb{a},\pmb{b},\pmb{c}\},S,
\{S\to S\pmb{a}\mid S\pmb{c}\mid A\pmb{b}\mid \lambda,\;
  A\to S\pmb{a}\mid S\pmb{c}\mid \lambda\}\rangle
\]
Here $S$ generates the strings with no consecutive $\pmb{b}$s, while $A$
generates those whose leftmost symbol is a single $\pmb{b}$.

\item Containing two consecutive $\pmb{b}$s:
\[
G = \langle\{S,A,B\},\{\pmb{a},\pmb{b},\pmb{c}\},S,
\{S\to S\pmb{a}\mid S\pmb{c}\mid A\pmb{b},\;
  A\to S\pmb{a}\mid S\pmb{c}\mid B\pmb{b},\;
  B\to B\pmb{a}\mid B\pmb{b}\mid B\pmb{c}\mid \lambda\}\rangle
\]

\item Even number of $\pmb{a}$s:
\[
G = \langle\{S,A\},\{\pmb{a},\pmb{b},\pmb{c}\},S,
\{S\to S\pmb{b}\mid S\pmb{c}\mid A\pmb{a}\mid \lambda,\;
  A\to S\pmb{a}\mid A\pmb{b}\mid A\pmb{c}\}\rangle
\]

\item Not ending in $\pmb{b}$:
\[
G = \langle\{S,A\},\{\pmb{a},\pmb{b},\pmb{c}\},S,
\{S\to A\pmb{a}\mid A\pmb{c}\mid \lambda,\;
  A\to A\pmb{a}\mid A\pmb{b}\mid A\pmb{c}\mid \lambda\}\rangle
\]

\item No $\pmb{c}$s:
\[
G = \langle\{S\},\{\pmb{a},\pmb{b},\pmb{c}\},S,
\{S\to S\pmb{a}\mid S\pmb{b}\mid \lambda\}\rangle
\]
\end{enumerate}

%%------------------------------------------------------------
\item \textbf{Complete the inductive portion of the proof of Theorem 8.8.}

Let $G=\langle\Omega,\Sigma,S,P\rangle$ already be in the normal form of
Theorem~8.7, and let $G^0=\langle\Omega\cup\{Z\},\Sigma,Z,P^0\rangle$ be the
grammar constructed in Theorem~8.8.

The needed induction is on the number of productions used after the initial
step $Z\to S$.

\textbf{Claim}: For every $A\in\Omega$ and every $x\in\Sigma^+$,
\[
A\DRightarrow_{G^0} x \quad\text{if and only if}\quad A\DRightarrow_G x.
\]

We prove the claim by induction on the number of productions in the derivation.

\textbf{Base}: A one-step derivation in $G^0$ must use a production
$A\to\pmb{a}$. By construction of $P^0$, there is some $B\in\Omega$ such that
$A\to\pmb{a}B\in P$ and $B\to\lambda\in P$. Hence in $G$ we have
$A\Rightarrow\pmb{a}B\Rightarrow\pmb{a}$. Conversely, if $A\DRightarrow_G x$ in
two steps, then necessarily $x=\pmb{a}$, with productions
$A\to\pmb{a}B,\ B\to\lambda$, so $A\to\pmb{a}$ is in $P^0$ and
$A\DRightarrow_{G^0}\pmb{a}$.

\textbf{Step}: Assume the claim holds for derivations using at most $n$
productions, and suppose $A\DRightarrow_{G^0} x$ uses $n+1$ productions.

If the first production is $A\to\pmb{a}B$, then
\[
A\Rightarrow_{G^0}\pmb{a}B\DRightarrow_{G^0} x,
\]
so $x=\pmb{a}y$ and $B\DRightarrow_{G^0} y$ uses $n$ productions. By the
induction hypothesis, $B\DRightarrow_G y$, and therefore
$A\Rightarrow_G\pmb{a}B\DRightarrow_G x$.

The other possibility is that the first production is $A\to\pmb{a}$; then
$x=\pmb{a}$ and this is already the base case.

For the converse direction, suppose $A\DRightarrow_G x$ and the derivation uses
$n+2$ productions. The first production cannot be $A\to\lambda$, since
$x\in\Sigma^+$. If the first production is $A\to\pmb{a}B$, and the remainder is
$B\DRightarrow_G y$ with $x=\pmb{a}y$, then either $y\in\Sigma^+$, in which
case the induction hypothesis gives $B\DRightarrow_{G^0} y$ and hence
$A\Rightarrow_{G^0}\pmb{a}B\DRightarrow_{G^0} x$, or else $y=\lambda$, in which
case $B\to\lambda\in P$ and therefore $A\to\pmb{a}\in P^0$, so
$A\DRightarrow_{G^0}\pmb{a}=x$.

Now apply the claim with $A=S$. Any terminal derivation in $G^0$ either is
$Z\Rightarrow\lambda$ (which occurs exactly when $S\to\lambda\in P$) or begins
$Z\Rightarrow S$ and then follows a derivation covered by the claim. Hence
$L(G^0)=L(G)$. \qed

%%------------------------------------------------------------
\item \textbf{Complete the inductive portion of the proof of Theorem 8.5.}

Let $G_*=\langle\Omega\cup\{Z\},\Sigma,Z,P_*\rangle$ be the grammar of
Theorem~8.5. The theorem says that derivations using exactly $n$ productions of
the form $A\to xZ$ generate exactly the words in $L(G)^n$.

\textbf{Claim}: For every $n\in\NN$ and every $w\in\Sigma^*$,
\[
Z\DRightarrow_{G_*} w \text{ using exactly $n$ productions of the form }
A\to xZ
\]
if and only if $w\in L(G)^n$.

We prove this by induction on $n$.

\textbf{Base} ($n=0$): The only terminal derivation from $Z$ that uses no
production of the form $A\to xZ$ is $Z\Rightarrow\lambda$. Hence the words
generated are exactly $L(G)^0=\{\lambda\}$.

\textbf{Step}: Assume the claim holds for $n$, and consider a derivation using
exactly $n+1$ productions of the form $A\to xZ$. Since the start symbol $Z$
appears only in the productions $Z\to\lambda$ and $Z\to S$, such a derivation
must begin
\[
Z\Rightarrow S\DRightarrow_{G_*} uZ\DRightarrow_{G_*} uv,
\]
where the portion $S\DRightarrow_{G_*} uZ$ uses exactly one production
$A\to xZ$. Replacing that final production by the corresponding production
$A\to x$ from $G$ shows that $u\in L(G)$. The remainder
$Z\DRightarrow_{G_*} v$ uses exactly $n$ productions of the form $A\to xZ$, so
by the induction hypothesis $v\in L(G)^n$. Therefore
$w=uv\in L(G)\cdot L(G)^n=L(G)^{n+1}$.

Conversely, let $w\in L(G)^{n+1}$. Then $w=uv$ for some
$u\in L(G)$ and $v\in L(G)^n$. Choose a derivation $S\DRightarrow_G u$ in $G$.
Its last step uses some production $A\to x$ in $P$; in $G_*$ this production
has been replaced by $A\to xZ$, so the same derivation yields
$S\DRightarrow_{G_*} uZ$. By the induction hypothesis,
$Z\DRightarrow_{G_*} v$ using exactly $n$ productions of the form $A\to xZ$.
Thus
\[
Z\Rightarrow S\DRightarrow_{G_*} uZ\DRightarrow_{G_*} uv=w
\]
uses exactly $n+1$ such productions. Hence the words generated in this way are
exactly the words in $L(G)^{n+1}$. \qed

%%------------------------------------------------------------
\item \textbf{Use the efficient algorithm of Theorem 8.3 to find regular expressions for $L(G_5)$, $L(G_6)$, $L(G_7)$ in Exercise 8.3.}

The efficient algorithm eliminates one variable at a time using
Theorem 6.1.

\begin{enumerate}
\item[$G_5$:] Equations: $S_0 = \pmb{0}S_2 \cup \pmb{1}S_1 \cup \lambda$, $S_1 = \pmb{0}S_1 \cup \pmb{1}S_2$, $S_2 = \pmb{0}S_2 \cup \pmb{1}S_0$.
Solve for $S_2$ from third equation: $S_2 = \pmb{0}^*\pmb{1}S_0$.
Substitute into $S_1$: $S_1 = \pmb{0}S_1 \cup \pmb{1}\pmb{0}^*\pmb{1}S_0$, so $S_1 = \pmb{0}^*\pmb{1}\pmb{0}^*\pmb{1}S_0$.
Substitute $S_1$ and $S_2$ into $S_0$: $S_0 = \pmb{0}\cdot\pmb{0}^*\pmb{1}S_0 \cup \pmb{1}\cdot\pmb{0}^*\pmb{1}\pmb{0}^*\pmb{1}S_0 \cup \lambda = (\pmb{0}^+\pmb{1} \cup \pmb{1}\pmb{0}^*\pmb{1}\pmb{0}^*\pmb{1})S_0 \cup \lambda$.
By Theorem 6.1,
\[
L(G_5) = (\pmb{0}^+\pmb{1}\cup
\pmb{1}\pmb{0}^*\pmb{1}\pmb{0}^*\pmb{1})^*.
\]

\item[$G_6$:] Equations: $Z = \pmb{a}Z \cup \pmb{b}T$, $T = \pmb{a}Z$. Substitute: $Z = \pmb{a}Z \cup \pmb{ba}Z = (\pmb{a}\cup\pmb{ba})Z \cup \emptyset$. By Theorem 6.1, $Z = (\pmb{a}\cup\pmb{ba})^*\cdot\emptyset = \emptyset$. So $L(G_6) = \emptyset$.

\item[$G_7$:] Equations: $S = \pmb{a}S \cup \pmb{ab}B \cup \pmb{c}C$, $B = \pmb{ab}B \cup \lambda$, $C = \pmb{c}C \cup \pmb{ca}$.
Solve $B = (\pmb{ab})^*$.
Solve $C = \pmb{c}^*\pmb{ca}$.
Substitute: $S = \pmb{a}S \cup \pmb{ab}(\pmb{ab})^* \cup \pmb{c}\cdot\pmb{c}^*\pmb{ca} = \pmb{a}S \cup (\pmb{ab})^+\cup\pmb{c}^+\pmb{a}$.
By Theorem 6.1:
\[
L(G_7) = \pmb{a}^*((\pmb{ab})^+\cup\pmb{c}^+\pmb{a})
       = \pmb{a}^*(\pmb{ab})^+\cup\pmb{a}^*\pmb{c}^+\pmb{a}.
\]
\end{enumerate}

%%------------------------------------------------------------
\item \textbf{Left-linear version of Theorem 8.3 and Lemmas 8.4--8.5.}

\begin{enumerate}
\item Let
\[
G=\langle\{S_1,S_2,\ldots,S_n\},\Sigma,S_1,P\rangle
\]
be a left-linear grammar, and for each nonterminal $S_i$ let $X_{S_i}$ be the
set of terminal strings derivable from $S_i$.  Then
\[
X_{S_k}=E_k\cup X_{S_1}A_{1k}\cup X_{S_2}A_{2k}\cup\cdots\cup X_{S_n}A_{nk},
\qquad k=1,2,\ldots,n,
\]
where $E_k$ is the union of all terminal strings $x$ that appear in
productions of the form $S_k\to x$, and $A_{jk}$ is the union of all terminal
strings $x$ that appear in productions of the form $S_k\to S_jx$.

\item The left-linear analogue of Lemma~8.4 is: if
\[
X=E\cup XA,
\]
then $EA^*$ is always a solution, and any other solution $Y$ must satisfy
$EA^*\subseteq Y$.

The left-linear analogue of Lemma~8.5 is obtained by eliminating the last
unknown on the right.  If
\[
X_n=E_n\cup X_1A_{1n}\cup\cdots\cup X_{n-1}A_{(n-1)n}\cup X_nA_{nn},
\]
then the compatible minimal choice for $X_n$ is
\[
X_n=(E_n\cup X_1A_{1n}\cup\cdots\cup X_{n-1}A_{(n-1)n})A_{nn}^*.
\]
Substituting this into the first $n-1$ equations yields
\[
\hat{E}_i=E_i\cup E_nA_{nn}^*A_{ni},
\qquad
\hat{A}_{ji}=A_{ji}\cup A_{jn}A_{nn}^*A_{ni},
\]
for $i,j=1,2,\ldots,n-1$.

\item For $G_2$ in Example~8.12, the language equations are
\[
X_S=X_A\pmb{baa},\qquad X_A=\lambda\cup X_A\pmb{a}.
\]
By the left-linear version of Theorem 6.1,
\[
X_A=\lambda\pmb{a}^*=\pmb{a}^*.
\]
Therefore
\[
L(G_2)=X_S=X_A\pmb{baa}=\pmb{a}^*\pmb{baa}.
\]
\end{enumerate}

%%------------------------------------------------------------
\item \textbf{Grammar $Q$ for binary numbers with decimal point.}

\begin{enumerate}
\item A correct NDFA is
\[
A=\langle\{\pmb{0},\pmb{1},\pmb{.}\},\{I,D,F\},I,\delta,\{F\}\rangle,
\]
where
\[
\delta(I,\pmb{0})=\{I,D\},\qquad \delta(I,\pmb{1})=\{I,D\},
\]
\[
\delta(D,\pmb{.})=\{F\},\qquad
\delta(F,\pmb{0})=\{F\},\qquad \delta(F,\pmb{1})=\{F\},
\]
and all other transitions are empty.  The state $D$ means ``the last digit
before the decimal point has just been read.''  This accepts strings such as
$\pmb{011.}$ and $\pmb{101.11}$, but not strings beginning with a dot.

\item Right-linear equations:
\begin{align*}
I &= \pmb{0}I \cup \pmb{1}I \cup \pmb{0}\pmb{.}F \cup \pmb{1}\pmb{.}F \\
F &= \lambda \cup \pmb{0}F \cup \pmb{1}F
\end{align*}

\item Solving: $F = (\pmb{0}\cup\pmb{1})^*$ by Theorem 6.1 applied to $F = (\pmb{0}\cup\pmb{1})F\cup\lambda$.
Substituting into $I$: $I = (\pmb{0}\cup\pmb{1})I \cup (\pmb{0}\cup\pmb{1})\pmb{.}(\pmb{0}\cup\pmb{1})^*$.
By Theorem 6.1, $I = (\pmb{0}\cup\pmb{1})^+\pmb{.}(\pmb{0}\cup\pmb{1})^*$.
\end{enumerate}

%%------------------------------------------------------------
\item \textbf{Right-linear grammars for the given regular expressions.}

\begin{enumerate}
\item $(\pmb{a}\cup\pmb{b})\pmb{c}^*(\pmb{d}\cup(\pmb{ab})^*)$:
\begin{align*}
G &= \langle\{S,A,B,C\},\{\pmb{a},\pmb{b},\pmb{c},\pmb{d}\},S,P\rangle \\
P &= \{S\to\pmb{a}A|\pmb{b}A,\ A\to\pmb{c}A|\pmb{d}|\pmb{a}B|\lambda,\ B\to\pmb{b}C,\ C\to\pmb{a}B|\lambda\}
\end{align*}
State $A$ handles $\pmb{c}^*$ and the choice of $\pmb{d}$ or $(\pmb{ab})^*$; state $B$ after seeing $\pmb{a}$ in the $(\pmb{ab})^*$ part; state $C$ after $\pmb{ab}$.

\item $(\pmb{a}\cup\pmb{b})^*\pmb{a}(\pmb{a}\cup\pmb{b})^*$ (strings containing at least one $\pmb{a}$):
\begin{align*}
G &= \langle\{S,A\},\{\pmb{a},\pmb{b}\},S,\{S\to\pmb{a}S|\pmb{b}S|\pmb{a}A,\ A\to\pmb{a}A|\pmb{b}A|\lambda\}\rangle
\end{align*}
State $S$ generates the arbitrary prefix before the required occurrence of
$\pmb{a}$; the production $S\to\pmb{a}A$ supplies that required occurrence.
State $A$ then generates the arbitrary suffix and may terminate by
$A\to\lambda$.
\end{enumerate}

%%------------------------------------------------------------
\item \textbf{Left-linear grammars for the same expressions.}

\begin{enumerate}
\item $(\pmb{a}\cup\pmb{b})\pmb{c}^*(\pmb{d}\cup(\pmb{ab})^*)$: Reverse the order of derivation.
\[
G = \langle\{S,A,B\},\{\pmb{a},\pmb{b},\pmb{c},\pmb{d}\},S,
\{S\to A\pmb{d}|B,\ A\to A\pmb{c}|\pmb{a}|\pmb{b},\ B\to C\pmb{b}|A,\ C\to B\pmb{a}\}\rangle
\]
Here $A$ generates $(\pmb{a}\cup\pmb{b})\pmb{c}^*$, so the production
$S\to A\pmb{d}$ gives the $\pmb{d}$ branch.  The other branch starts with
$S\to B$, then $B\to A$ gives the $(\pmb{ab})^0$ case, while the cycle
$B\to C\pmb{b}\to B\pmb{ab}$ adds one more copy of $(\pmb{ab})$ on the right.

\item $(\pmb{a}\cup\pmb{b})^*\pmb{a}(\pmb{a}\cup\pmb{b})^*$:
\[
G = \langle\{S,A\},\{\pmb{a},\pmb{b}\},S,\{S\to S\pmb{a}|S\pmb{b}|A\pmb{a},\ A\to A\pmb{a}|A\pmb{b}|\lambda\}\rangle
\]
Here $A$ generates the arbitrary prefix, and the production $S\to A\pmb{a}$
places the required $\pmb{a}$ after that prefix.  The productions
$S\to S\pmb{a}$ and $S\to S\pmb{b}$ append the arbitrary suffix on the right.
\end{enumerate}

%%------------------------------------------------------------
\item \textbf{Algorithm to convert right-linear to left-linear grammar.}

\begin{enumerate}
\item \textbf{Algorithm}:
\begin{enumerate}
\item From the right-linear grammar $G$, construct the NDFA $A_G$ of
Lemma~8.2.

\item Reverse all arrows in $A_G$ and interchange initial and final states.  As
in the reverse construction used in Exercise~5.20 and cited in the proof of
Theorem 8.2, the resulting machine accepts $L(G)^r$.

\item Apply the subset construction from Definition 4.5 to obtain a
DFA for $L(G)^r$.

\item Apply Lemma~8.1 to that DFA to get a right-linear grammar $H$ for
$L(G)^r$.

\item Reverse every production in $H$.  By Lemma~8.3, the resulting grammar is
left-linear and generates
\[
L(H)^r=(L(G)^r)^r=L(G).
\]
\end{enumerate}

\item Applying this to $G_4$: first construct $A_{G_4}$ from Lemma~8.2, then
reverse it, determinize, and apply Lemma~8.1.  The resulting right-linear
grammar for $L(G_4)^r$ is
\[
H=\langle\{R,U,V,W,X,L,M,N,P\},\Sigma,R,P_H\rangle,
\]
where $P_H$ contains
\[
\begin{aligned}
R&\to \pmb{b}U\\
U&\to \pmb{b}V\mid \pmb{a}W\\
V&\to \pmb{a}L\\
W&\to \pmb{b}X\\
X&\to \pmb{b}L\\
L&\to \lambda\mid \pmb{c}M\\
M&\to \pmb{b}N\\
N&\to \pmb{b}P\\
P&\to \pmb{a}L.
\end{aligned}
\]
Reversing every production now yields a left-linear grammar for $L(G_4)$:
\[
\begin{aligned}
R&\to U\pmb{b}\\
U&\to V\pmb{b}\mid W\pmb{a}\\
V&\to L\pmb{a}\\
W&\to X\pmb{b}\\
X&\to L\pmb{b}\\
L&\to \lambda\mid M\pmb{c}\\
M&\to N\pmb{b}\\
N&\to P\pmb{b}\\
P&\to L\pmb{a}.
\end{aligned}
\]
\end{enumerate}

%%------------------------------------------------------------
\item \textbf{Algorithm to convert a regular grammar to one generating the complement.}

\begin{enumerate}
\item Construct an automaton for $L(G)$ from grammar $G$ using Lemma~8.2, and
then, if necessary, convert it to an equivalent DFA $A$ by Definition 4.5.
\item Compute the complement DFA $A^c$ by keeping the same states, start state,
and transition function as $A$, and replacing the final-state set $F$ by
$S-F$.  Since $A$ is a DFA, its transition function is already total; the
earlier parenthetical ``after making $A$ complete/total'' was only meant as a
guard against stopping with a partial transition table before the DFA step has
actually been completed.
\item Construct a right-linear grammar $G^c$ from $A^c$.
Then $L(G^c) = \overline{L(G)}$.
\end{enumerate}

%%------------------------------------------------------------
\item \textbf{Algorithm to determine whether $L(G) = \emptyset$ for a regular grammar $G$.}

\begin{enumerate}
\item Build the automaton $A_G$ from $G$.
\item Compute the connected state set $S^c$ from the start state by the
successive $C_i$ construction of Definition 3.10.  The assertion
proved in Exercise~3.53 justifies that these $C_i$ sets produce exactly the
states reachable by strings of bounded length, so the iteration is the
textbook reachability algorithm.
\item If no final state is reachable, $L(G) = \emptyset$; otherwise $L(G) \neq \emptyset$.
\end{enumerate}
This terminates since the automaton has finitely many states.

%%------------------------------------------------------------
\item \textbf{Algorithm to determine whether $L(G)$ is infinite for a regular grammar $G$.}

\begin{enumerate}
\item Build the automaton $A_G$ from $G$.
\item Find all states reachable from the start state.
\item Find all states from which a final state is reachable.
\item Check whether the intersection of these two sets (the ``useful'' states) contains any cycle (e.g., by DFS looking for back edges).
\item $L(G)$ is infinite iff a cycle exists among the useful states.
\end{enumerate}

%%------------------------------------------------------------
\item \textbf{Algorithm to determine whether two right-linear grammars $G_1$ and $G_2$ generate the same language.}

\begin{enumerate}
\item Construct DFAs
$A_1=\langle\Sigma,S_1,s_{0_1},\delta_1,F_1\rangle$ and
$A_2=\langle\Sigma,S_2,s_{0_2},\delta_2,F_2\rangle$ from $G_1$ and $G_2$,
using Lemma~8.2 followed by Definition 4.5 when determinization is
needed.  Rename states if necessary so that $S_1\cap S_2=\emptyset$.
\item Form the composite DFA
\[
A=\langle\Sigma,S_1\cup S_2,s_{0_1},\delta,F_1\cup F_2\rangle,
\]
where $\delta(s,\pmb{a})=\delta_1(s,\pmb{a})$ for $s\in S_1$ and
$\delta(s,\pmb{a})=\delta_2(s,\pmb{a})$ for $s\in S_2$.
\item Use the Chapter~3 algorithm for computing $E_A$ (Corollary 3.5).
\item Answer yes exactly when $s_{0_1}E_As_{0_2}$.  If the two start states are
state-equivalent in the composite machine, every input word either reaches final
states in both component machines or reaches nonfinal states in both; if they
are not equivalent, some word is accepted by exactly one component machine.
\end{enumerate}

%%------------------------------------------------------------
\item \textbf{Determine $L(H)$ for grammar $H$.}

The grammar is $H = \langle\{A,B,S\},\{\pmb{a},\pmb{b},\pmb{c}\},S,\{S\to\pmb{a}SB\pmb{c},S\to\lambda,SB\to\pmb{b}S,\pmb{c}B\to B\pmb{c}\}\rangle$.

This grammar generates $\lambda$, and for every $n\ge 2$ it generates
$\pmb{a}^n\pmb{b}^n\pmb{c}^n$.  For example,
\[
S\Rightarrow \pmb{aa}SB\pmb{c}B\pmb{c}
\Rightarrow \pmb{aa}SBB\pmb{cc}
\Rightarrow \pmb{aab}SB\pmb{cc}
\Rightarrow \pmb{aabb}S\pmb{cc}
\Rightarrow \pmb{aabbcc}.
\]
In general, after $n$ uses of $S\to\pmb{a}SB\pmb{c}$ we have
\[
\pmb{a}^nSB^n\pmb{c}^n.
\]
Repeatedly use $\pmb{c}B\to B\pmb{c}$ to move one $B$ left until it is adjacent
to $S$, then apply $SB\to\pmb{b}S$.  Doing this $n$ times yields
\[
\pmb{a}^n\pmb{b}^nS\pmb{c}^n,
\]
and then $S\to\lambda$ gives $\pmb{a}^n\pmb{b}^n\pmb{c}^n$.

However, $\pmb{abc}$ is \emph{not} generated: after one use of
$S\to\pmb{a}SB\pmb{c}$ and then $S\to\lambda$ we obtain $\pmb{a}B\pmb{c}$, and
there is no rule that can eliminate that lone $B$.  Thus the nonempty terminal
strings begin at length $6$.

Therefore
\[
L(H)=\{\lambda\}\cup\{\pmb{a}^n\pmb{b}^n\pmb{c}^n\mid n\ge 2\}.
\]

%%------------------------------------------------------------
\item \textbf{What is wrong with the concatenation construction $G^\bullet = \langle\Omega_1\cup\Omega_2\cup\{Z\},\Sigma,Z,P_1\cup P_2\cup\{Z\to S_1\cdot S_2\}\rangle$?}

The production $Z \to S_1 S_2$ has two nonterminals on the right-hand side. In a right-linear grammar, all productions must have the form $A \to w$ or $A \to wB$ (a string of terminals followed by at most one nonterminal). The production $S_1 S_2$ violates this since $S_1$ is a nonterminal followed by $S_2$ (also nonterminal), with no terminal prefix. This is a context-free production, not a right-linear one, so the resulting grammar is not right-linear and cannot directly represent a language in \GSIG.

%%------------------------------------------------------------
\item \textbf{Prove \GSIG\ is closed under concatenation.}

\begin{enumerate}
\item Construction: Let $G_1 = \langle\Omega_1,\Sigma,S_1,P_1\rangle$ and $G_2 = \langle\Omega_2,\Sigma,S_2,P_2\rangle$ be right-linear grammars with $\Omega_1 \cap \Omega_2 = \emptyset$. Define:
\[
G^\bullet = \langle\Omega_1\cup\Omega_2,\Sigma,S_1,P^\bullet\rangle
\]
where $P^\bullet$ is obtained from $P_1$ by replacing each terminal production $A \to w$ (where $A \in \Omega_1$) with $A \to wS_2$, plus all productions from $P_2$ unchanged.

\item Proof that $L(G^\bullet) = L(G_1)\cdot L(G_2)$:
If $u \in L(G_1)$ and $v \in L(G_2)$, then $S_1 \DRightarrow u$ in $G_1$ via some derivation ending with $A \to u'$ (terminal production). In $G^\bullet$, this becomes $S_1 \DRightarrow u'S_2 \DRightarrow u'v = uv$.
Conversely, any derivation in $G^\bullet$ must use rules from $P_1$ until reaching a production $A \to wS_2$, producing some prefix, then switch to $P_2$. \qed
\end{enumerate}

%%------------------------------------------------------------
\item \textbf{Construct a $\lambda$-NDFA $A'$ with one start and one final state from an NDFA $A$ without $\lambda$-transitions.}

Let $A = \langle\Sigma,S,s_0,\delta,F\rangle$ be an NDFA without $\lambda$-transitions.

Add two new states $i$ and $f$, where $i$ will be the unique start state and
$f$ the unique final state.  Define
\[
A'=\langle\Sigma,S\cup\{i,f\},i,\delta',\{f\}\rangle,
\]
where
\[
\delta'(s,\pmb{a})=\delta(s,\pmb{a})\qquad(s\in S,\ \pmb{a}\in\Sigma),
\]
\[
\delta'(i,\lambda)=\{s_0\},
\]
and for every old final state $p\in F$ we add the $\lambda$-transition
\[
f\in\delta'(p,\lambda).
\]
There are no other new transitions.

Every accepting path of $A$ becomes an accepting path of $A'$ by first taking
$i\to s_0$, following the old transitions of $A$, and then taking a final
$\lambda$-move into $f$.  Conversely, any accepting path in $A'$ must begin with
$i\to s_0$, use only old transitions while reading the input, and end with a
$\lambda$-move from an old final state into $f$.  Hence
\[
L(A')=L(A),
\]
and $A'$ has exactly one start state and exactly one final state.

%%------------------------------------------------------------
\item \textbf{Complete the proof of Lemma 8.2.}

\begin{enumerate}
\item Inductive statement: For all $n \geq 0$ and all $x \in \Sigma^n$ and all nonterminals $A$: $A \DRightarrow x$ in $G$ iff $\DBAR_A(A, x) \in F$.

\item Proof by induction on $n$:
\textbf{Base} ($n=0$): $A \DRightarrow \lambda$ iff $A \to \lambda \in P$ iff $A \in F$ in the NDFA iff $\DBAR(A,\lambda) = A \in F$.
\textbf{Step}: Assume the result holds for length $n$. For $x = ay$ ($|y|=n$): $A \DRightarrow ay$ iff $(\exists B)$ $A \to \pmb{a}B \in P$ and $B \DRightarrow y$ iff $(\exists B)$ $B \in \delta(A,\pmb{a})$ and (by IH) $\DBAR(B,y) \in F$ iff $\DBAR(A,ay) \in F$. \qed
\end{enumerate}

%%------------------------------------------------------------
\item \textbf{Complete the proof of Lemma 8.3.}

\begin{enumerate}
\item Inductive statement: for every $n\ge 0$, every pair of nonterminals
$A,B$, and every $x\in\Sigma^*$,
\[
A\DRightarrow_G^{ n}Bx
\qquad\Longleftrightarrow\qquad
A\DRightarrow_{G^r}^{ n}x^rB.
\]

\item Proof by induction on $n$.

\textbf{Base} ($n=0$): A zero-step derivation has the form $A\DRightarrow^0 A$,
so $B=A$ and $x=\lambda$.  Then also
\[
A\DRightarrow_{G^r}^0\lambda^rA=A.
\]

\textbf{Step}: Assume the statement true for $n$ productions.  Suppose
\[
A\DRightarrow_G^{ n+1}By.
\]
Then the first production used in $G$ has the form
\[
A\to Cz
\]
for some $C\in\Omega\cup\{\lambda\}$ and some $z\in\Sigma^*$, and after that
\[
C\DRightarrow_G^{ n}Bx
\]
with $y=zx$.  By the induction hypothesis,
\[
C\DRightarrow_{G^r}^{ n}x^rB.
\]
Since $G^r$ contains the reversed production
\[
A\to z^rC,
\]
we obtain
\[
A\Rightarrow_{G^r} z^rC \DRightarrow_{G^r}^{ n} z^r x^r B = (xz)^r B = y^r B.
\]
The converse implication is identical with the roles of $G$ and $G^r$
reversed.  Therefore the statement holds for all $n$.  In particular,
terminal derivations are reversed, so
\[
L(G^r)=L(G)^r.
\]
\qed
\end{enumerate}

%%------------------------------------------------------------
\item \textbf{Reasons for the assertions in the second half of the proof of Theorem 8.2.}

The second half of Theorem~8.2 begins with a language $L$ generated by a
right-linear grammar.

\begin{enumerate}
\item Since $L$ is generated by a right-linear grammar, Theorem~8.1 implies
that $L\in\DSIG$.

\item Exercise~5.20 (or 6.36) shows that \DSIG\ is closed under reverse, so
$L^r\in\DSIG$.

\item Because $\DSIG=\GSIG$, there is a right-linear grammar $G$ such that
$L(G)=L^r$.

\item Lemma~8.3 says that $G^r$ is a left-linear grammar and that
\[
L(G^r)=L(G)^r.
\]
Since $L(G)=L^r$, this becomes
\[
L(G^r)=(L^r)^r=L.
\]

\item Therefore $G^r$ is a left-linear grammar that generates $L$.  This is the
``corresponding left-linear grammar'' claimed in the theorem.
\end{enumerate}

%%------------------------------------------------------------
\item \textbf{Formal inductive proofs.}

\begin{enumerate}
\item $L(G_1) = \{\pmb{a}^n\pmb{baa} \mid n \in \NN\}$ (Example 8.7):

\textbf{Claim}: For all $n \geq 0$, $S \DRightarrow \pmb{a}^n\pmb{baa}$ in $G_1$.
\textbf{Base} ($n=0$): $S \Rightarrow \pmb{baa}$ (production $S \to \pmb{baa}$).
\textbf{Step}: If $S \DRightarrow \pmb{a}^n\pmb{baa}$, then $S \Rightarrow \pmb{a}S \DRightarrow \pmb{a}\cdot\pmb{a}^n\pmb{baa} = \pmb{a}^{n+1}\pmb{baa}$.
Conversely, any derivation from $S$ either uses $S \to \pmb{baa}$ (giving $\lambda\cdot\pmb{baa}$) or $S \to \pmb{a}S$ (giving a longer string by IH). So $L(G_1) = \{\pmb{a}^n\pmb{baa}\}$. \qed

\item $L(G_8) = \{\pmb{a}^n\pmb{baa} \mid n \in \NN\}$ (Example 8.9):
\textbf{Claim}: For every $n \geq 0$, $Z \DRightarrow \pmb{a}^n\pmb{baa}$ in
$G_8$.

\textbf{Base} ($n=0$): $Z \Rightarrow \pmb{baa}$.

\textbf{Step}: If $Z \DRightarrow \pmb{a}^n\pmb{baa}$, then
$Z \Rightarrow \pmb{a}Z \DRightarrow \pmb{a}^{n+1}\pmb{baa}$.

Conversely, any derivation from $Z$ either uses $Z\to\pmb{baa}$ immediately or
begins with $Z\to\pmb{a}Z$. Repeating the second choice $n$ times and then
using $Z\to\pmb{baa}$ yields exactly the string $\pmb{a}^n\pmb{baa}$, and there
is no other way to terminate a derivation. Therefore
$L(G_8)=\{\pmb{a}^n\pmb{baa}\mid n\in\NN\}$. \qed
\end{enumerate}

%%------------------------------------------------------------
\item \textbf{Mixed grammar $C$ with both right-linear and left-linear productions.}

$C = \langle\{S,A,B\},\{\pmb{0},\pmb{1}\},S,\{S\to\pmb{0}A|\pmb{1}B|\pmb{0}|\pmb{1}|\lambda,A\to S\pmb{0},B\to S\pmb{1}\}\rangle$.

\begin{enumerate}
\item $L(C) = \{w \in \{\pmb{0},\pmb{1}\}^* \mid w = w^r\}$ --- the set of all palindromes over $\{\pmb{0},\pmb{1}\}$.

Derivation: $S \Rightarrow \pmb{0}A \Rightarrow \pmb{0}S\pmb{0}$. Each step wraps the current sentential form between matching symbols. The grammar generates all even-length and odd-length palindromes.

\item Yes, $L(C) = \{w \mid w = w^r\}$ is not regular (it requires unbounded memory to match the first and last symbols). So $L(C)$ is FAD only for the trivial case; in general, $L(C)$ is not FAD.

\item The definition of regular grammars should \emph{not} be expanded to include such mixed grammars, because the resulting class of languages would no longer coincide with the regular languages. Mixed grammars generate languages (like palindromes) that are context-free but not regular.
\end{enumerate}

%%------------------------------------------------------------
\item \textbf{Importance of $\Omega_1 \cap \Omega_2 = \emptyset$ in Theorem 8.4.}

\begin{enumerate}
\item \textbf{Why important}: If $\Omega_1 \cap \Omega_2 \neq \emptyset$, then a shared nonterminal $A$ would have its productions from both $P_1$ and $P_2$ merged. Derivations starting from $S_1$ could use $P_2$ productions for $A$ and vice versa, producing strings not in $L(G_1) \cup L(G_2)$.
\textbf{Example}: Let $G_1$ have $\{S_1 \to \pmb{a}A, A \to \pmb{b}\}$ and $G_2$ have $\{S_2 \to \pmb{c}A, A \to \pmb{d}\}$. If $A$ is shared, the union grammar with $Z \to S_1 | S_2$ would derive $\pmb{a}A$ and then use $A \to \pmb{d}$ from $G_2$, producing $\pmb{ad} \notin L(G_1) \cup L(G_2)$.

\item \textbf{Why possible}: We can always rename nonterminals in $G_2$ (choosing fresh names) to make $\Omega_1 \cap \Omega_2 = \emptyset$ without changing the language $L(G_2)$. This is a standard renaming/relabeling argument.
\end{enumerate}

%%------------------------------------------------------------
\item \textbf{If $A_G$ is disconnected, what does this say about $G$?}

If $A_G$ (the NDFA constructed from $G$) is disconnected, it means some states (nonterminals) are not reachable from the start state. This means the corresponding nonterminals in $G$ are \emph{useless}: they cannot appear in any derivation from the start symbol. The portions of $G$ defining those nonterminals contribute nothing to $L(G)$ and can be removed without affecting the language generated.

%%------------------------------------------------------------
\item \textbf{Apply Lemma 8.1 to the automata in Figure 8.4.}

Applying Lemma~8.1 means: use one nonterminal per state, use the start state as
start symbol, and add one production for each labeled transition; final states
also get $\lambda$-productions.

\begin{enumerate}
\item Figure~8.4(a):
\[
\begin{aligned}
r_0&\to \pmb{a}r_3\mid \pmb{b}r_1\mid \lambda\\
r_1&\to \pmb{a}r_0\mid \pmb{b}r_2\\
r_2&\to \pmb{a}r_1\mid \pmb{b}r_3\\
r_3&\to \pmb{a}r_2\mid \pmb{b}r_0.
\end{aligned}
\]

\item Figure~8.4(b):
\[
\begin{aligned}
q_0&\to \pmb{a}q_1\mid \pmb{b}q_0\mid \pmb{c}q_0\mid \lambda\\
q_1&\to \pmb{a}q_0\mid \pmb{b}q_0\mid \pmb{c}q_1.
\end{aligned}
\]

\item Figure~8.4(c):
\[
\begin{aligned}
S_0&\to \pmb{a}S_1\mid \pmb{b}S_2\\
S_1&\to \pmb{a}S_3\mid \pmb{b}S_2\\
S_2&\to \pmb{a}S_2\mid \pmb{b}S_2\mid \lambda\\
S_3&\to \pmb{a}S_0\mid \pmb{b}S_1.
\end{aligned}
\]

\item Figure~8.4(d):
\[
\begin{aligned}
t_0&\to \pmb{a}t_1\mid \pmb{b}t_2\\
t_1&\to \pmb{a}t_0\mid \pmb{b}t_2\mid \lambda\\
t_2&\to \pmb{a}t_3\mid \pmb{b}t_3\\
t_3&\to \pmb{a}t_1\mid \pmb{b}t_1\mid \lambda.
\end{aligned}
\]

\item Figure~8.4(e):
\[
\begin{aligned}
S&\to \pmb{a}r\mid \pmb{b}r\mid \pmb{b}t\mid \pmb{c}q\mid \lambda\\
r&\to \pmb{a}S\mid \pmb{b}t\mid \pmb{c}q\\
q&\to \pmb{a}r\mid \pmb{b}r\mid \pmb{c}q\mid \lambda\\
t&\to \pmb{a}r\mid \pmb{c}q.
\end{aligned}
\]
\end{enumerate}

%%------------------------------------------------------------
\item \textbf{Restate Lemma 8.1 for NDFAs.}

\begin{enumerate}
\item \textbf{Restated Lemma}: Given an NDFA
\[
A=\langle\Sigma,S,s_0,\delta,F\rangle,
\]
define
\[
G_A=\langle S,\Sigma,s_0,P\rangle
\]
by adding
\begin{itemize}
\item a production $s\to \pmb{a}t$ for every $t\in\delta(s,\pmb{a})$,
\item a production $s\to t$ for every $t\in\delta(s,\lambda)$,
\item a production $s\to\lambda$ for every final state $s\in F$.
\end{itemize}
Then $L(G_A)=L(A)$.

\item \textbf{Proof}: Every production records one NDFA move: letter transitions
become productions of the form $s\to \pmb{a}t$, $\lambda$-transitions become unit
productions $s\to t$, and final states get the production $s\to\lambda$.  Thus a
derivation in $G_A$ corresponds exactly to a path in $A$, and the derivation
ends in a terminal word precisely when the corresponding path ends in a final
state.  Therefore $L(G_A)=L(A)$.

\item Applied to Figure~8.5:
\begin{enumerate}
\item Figure~8.5(a):
\[
t_0\to \pmb{a}t_0\mid \pmb{c}t_0\mid \lambda.
\]

\item Figure~8.5(b):
\[
\begin{aligned}
q_1&\to \pmb{a}q_0\mid \pmb{b}q_1\mid \lambda\\
q_0&\to \pmb{b}q_1\mid \pmb{c}q_0.
\end{aligned}
\]

\item Figure~8.5(c):
\[
\begin{aligned}
r_0&\to \pmb{a}r_0\mid \pmb{a}r_1\mid \pmb{c}r_1\mid \lambda\\
r_1&\to \pmb{b}r_0.
\end{aligned}
\]

\item Figure~8.5(d):
\[
\begin{aligned}
S_0&\to \pmb{a}S_1\\
S_1&\to \pmb{a}S_2\mid \pmb{b}S_3\\
S_2&\to \lambda\\
S_3&\to \pmb{b}S_1.
\end{aligned}
\]

\item Figure~8.5(e):
\[
\begin{aligned}
S_1&\to \pmb{a}S_1\mid \pmb{a}S_2\\
S_2&\to \pmb{a}S_1\mid \pmb{b}S_1\mid \lambda.
\end{aligned}
\]

\item Figure~8.5(f):
\[
\begin{aligned}
S_0&\to \pmb{a}S_1\\
S_1&\to S_0\mid \pmb{b}S_2\\
S_2&\to \pmb{c}S_3\\
S_3&\to S_1\mid \lambda.
\end{aligned}
\]

\item Figure~8.5(g):
\[
\begin{aligned}
S_0&\to \pmb{a}S_1\mid \lambda\\
S_1&\to \pmb{a}S_0\mid \pmb{b}S_3\mid \pmb{c}S_2\\
S_2&\to \pmb{c}S_1\\
S_3&\to \pmb{b}S_1.
\end{aligned}
\]
\end{enumerate}
\end{enumerate}

%%------------------------------------------------------------
\item \textbf{Context-free grammars for languages over $\Sigma=\{\pmb{a},\pmb{b}\}$.}

\begin{enumerate}
\item $L_1$ = words where last letter matches first letter:
\[
G_1: S \to \pmb{a}X\pmb{a} \mid \pmb{b}X\pmb{b} \mid \pmb{a} \mid \pmb{b}, \quad X \to \pmb{a}X \mid \pmb{b}X \mid \lambda
\]
The empty word is excluded here, since it has no first or last letter.

\item $L_2$ = odd-length words where first letter matches center letter: Let $n=2k+1$.
\[
G_2: S \to \pmb{a} \mid \pmb{b} \mid \pmb{a}A\pmb{a} \mid \pmb{a}A\pmb{b} \mid \pmb{b}B\pmb{a} \mid \pmb{b}B\pmb{b},
\]
\[
A \to \pmb{a} \mid \pmb{a}A\pmb{a} \mid \pmb{a}A\pmb{b} \mid \pmb{b}A\pmb{a} \mid \pmb{b}A\pmb{b},
\]
\[
B \to \pmb{b} \mid \pmb{a}B\pmb{a} \mid \pmb{a}B\pmb{b} \mid \pmb{b}B\pmb{a} \mid \pmb{b}B\pmb{b}
\]
The nonterminal $A$ generates all odd-length words whose center letter is
$\pmb{a}$, and $B$ does the same for center letter $\pmb{b}$.  The productions
for $S$ force the first letter to match that center letter.

\item $L_3$ = words where last letter matches none of the other letters:
\[
G_3: S \to X\pmb{a} \mid X\pmb{b}, \quad \text{where } X \text{ generates strings not ending in } \pmb{a} \text{ (resp.\ } \pmb{b}\text{)}
\]
Since $\Sigma=\{\pmb{a},\pmb{b}\}$, ``last letter matches none of the others'' means all other letters differ. This is regular: last letter $\pmb{a}$, all others $\pmb{b}$; or last letter $\pmb{b}$, all others $\pmb{a}$: $L_3 = \pmb{b}^*\pmb{a} \cup \pmb{a}^*\pmb{b}$. So $G_3: S \to B\pmb{a} \mid A\pmb{b} \mid \pmb{a} \mid \pmb{b}$, $B \to B\pmb{b}\mid\lambda$, $A \to A\pmb{a}\mid\lambda$.

\item $L_4$ = even-length words where two center letters match:
\[
G_4: S \to \pmb{a}S\pmb{a} \mid \pmb{a}S\pmb{b} \mid \pmb{b}S\pmb{a} \mid \pmb{b}S\pmb{b} \mid \pmb{aa} \mid \pmb{bb}
\]

\item $L_5$ = odd-length words where center letter matches none of the others:
Since $\Sigma=\{\pmb{a},\pmb{b}\}$, if the center letter is $\pmb{a}$ then every
other letter must be $\pmb{b}$, and if the center letter is $\pmb{b}$ then every
other letter must be $\pmb{a}$. Therefore
\[
L_5=\{\pmb{b}^k\pmb{a}\pmb{b}^k\mid k\in\NN\}\cup
\{\pmb{a}^k\pmb{b}\pmb{a}^k\mid k\in\NN\},
\]
which is context-free but not regular:
\[
G_5: S \to A \mid B,\qquad A \to \pmb{b}A\pmb{b} \mid \pmb{a},\qquad B \to \pmb{a}B\pmb{a} \mid \pmb{b}
\]

\item Regular languages: $L_1$ is regular; $L_2$ is not regular; $L_3$ is regular ($\pmb{b}^*\pmb{a}\cup\pmb{a}^*\pmb{b}$); $L_4$ is not regular; $L_5$ is not regular.
\end{enumerate}

%%------------------------------------------------------------
\item \textbf{Context-free grammars for the given languages.}

\begin{enumerate}
\item $L = \{x \in \{\pmb{a},\pmb{b}\}^* \mid |x|_{\pmb{a}} < |x|_{\pmb{b}}\}$:
\[
G: S \to E\pmb{b}E \mid E\pmb{b}S,\qquad E \to \pmb{a}E\pmb{b}E \mid \pmb{b}E\pmb{a}E \mid \lambda
\]
Here $E$ generates exactly the strings with equally many $\pmb{a}$'s and
$\pmb{b}$'s.  The nonterminal $S$ inserts one or more extra $\pmb{b}$'s between
balanced pieces, so every derived string has strictly more $\pmb{b}$'s than
$\pmb{a}$'s, and every such string can be decomposed in that way.

\item $G = \{x \in \{\pmb{a},\pmb{b}\}^* \mid |x|_{\pmb{a}} \geq |x|_{\pmb{b}}\}$:
\[
\begin{aligned}
S &\to E \mid A\\
E &\to \pmb{a}E\pmb{b}E \mid \pmb{b}E\pmb{a}E \mid \lambda\\
A &\to E\pmb{a}E \mid E\pmb{a}A.
\end{aligned}
\]
Again $E$ generates the balanced strings.  The nonterminal $A$ inserts one or
more extra $\pmb{a}$'s between balanced pieces, so it generates exactly the
strings with more $\pmb{a}$'s than $\pmb{b}$'s.  Therefore $S$ generates
exactly the strings with $|x|_{\pmb{a}} \ge |x|_{\pmb{b}}$.

\item $K = \{w \in \{\pmb{0},\pmb{1}\}^* \mid w = w^r\}$ (palindromes):
$S \to \pmb{0}S\pmb{0} \mid \pmb{1}S\pmb{1} \mid \pmb{0} \mid \pmb{1} \mid \lambda$.

\item $\Phi = \{\pmb{a}^j\pmb{b}^k\pmb{c}^m \mid j\geq 3, k=m\}$:
$S \to \pmb{aaa}AT$, $A \to \pmb{a}A \mid \lambda$, $T \to \pmb{b}T\pmb{c} \mid \lambda$.
\end{enumerate}

%%------------------------------------------------------------
\item \textbf{More context-free grammars.}

\begin{enumerate}
\item $L_1 = \{x \in \{\pmb{a},\pmb{b}\}^* \mid |x|_{\pmb{a}} = 2|x|_{\pmb{b}}\}$:
\[
S \to \lambda
 \mid S\pmb{a}S\pmb{a}S\pmb{b}S
 \mid S\pmb{a}S\pmb{b}S\pmb{a}S
 \mid S\pmb{b}S\pmb{a}S\pmb{a}S
\]
Each nonempty production inserts exactly two $\pmb{a}$s and one $\pmb{b}$, with
recursive pieces before, between, and after them, so all generated strings
satisfy $|x|_{\pmb{a}} = 2|x|_{\pmb{b}}$.

\item $L_2 = \{x \in \{\pmb{a},\pmb{b}\}^* \mid |x|_{\pmb{a}} \neq |x|_{\pmb{b}}\}$:
\[
\begin{aligned}
S &\to A \mid B\\
E &\to \pmb{a}E\pmb{b}E \mid \pmb{b}E\pmb{a}E \mid \lambda\\
A &\to E\pmb{a}E \mid E\pmb{a}A\\
B &\to E\pmb{b}E \mid E\pmb{b}B.
\end{aligned}
\]
Here $E$ generates the words with equally many $\pmb{a}$'s and $\pmb{b}$'s.
Then $A$ inserts one or more extra $\pmb{a}$'s between balanced pieces, so it
generates exactly the words with more $\pmb{a}$'s than $\pmb{b}$'s. The
productions for $B$ do the mirror-image job with extra $\pmb{b}$'s, so $B$
generates exactly the words with more $\pmb{b}$'s than $\pmb{a}$'s.

\item Postfix expressions over $\{\pmb{A},\pmb{B},\pmb{+},\pmb{-}\}$:
$S \to \pmb{A} \mid \pmb{B} \mid SS\pmb{+} \mid SS\pmb{-}$.

\item Parenthesized infix expressions over $\{\pmb{A},\pmb{B},\pmb{+},\pmb{-},\pmb{(},\pmb{)}\}$:
$S \to \pmb{A} \mid \pmb{B} \mid \pmb{(}S\pmb{+}S\pmb{)} \mid \pmb{(}S\pmb{-}S\pmb{)}$.
\end{enumerate}

%%------------------------------------------------------------
\item \textbf{Context-sensitive grammars for the given languages.}

\begin{enumerate}
\item $\Gamma = \{w\pmb{2}w \mid w \in \{\pmb{0},\pmb{1}\}^*\}$:
A pure context-sensitive grammar for $\Gamma$ is
\[
G_\Gamma=\langle\{S,A,B\},\{\pmb{0},\pmb{1},\pmb{2}\},S,P\rangle,
\]
where
\begin{align*}
P=\{&S\to\pmb{2},\ S\to\pmb{0}A S,\ S\to\pmb{1}B S,\\
&A\pmb{0}\to\pmb{0}A,\ A\pmb{1}\to\pmb{1}A,\ B\pmb{0}\to\pmb{0}B,\ B\pmb{1}\to\pmb{1}B,\\
&A\pmb{2}\to\pmb{20},\ B\pmb{2}\to\pmb{21}\}.
\end{align*}

The productions $S\to\pmb{0}AS$ and $S\to\pmb{1}BS$ generate a sentential form
$M\pmb{2}w$, where the marker string $M$ records the symbols of $w$ in the same
left-to-right order.  The productions $A\pmb{2}\to\pmb{20}$ and
$B\pmb{2}\to\pmb{21}$ convert the rightmost remaining marker into the next
symbol of the left copy, while the swapping rules move the other markers right
until they too meet the $\pmb{2}$.  For example,
\[
S\Rightarrow \pmb{0}AS \Rightarrow \pmb{01}ABS \Rightarrow \pmb{01}AB\pmb{2}
\Rightarrow \pmb{01}A\pmb{21} \Rightarrow \pmb{0}A\pmb{121}
\Rightarrow \pmb{01201}.
\]
Thus $L(G_\Gamma)=\Gamma$.

\item $\Phi = \{\pmb{b}^{2^j} \mid j \in \NN\} = \{\pmb{b},\pmb{bb},\pmb{bbbb},\ldots\}$:
A pure context-sensitive grammar for $\Phi$ is
\[
G_\Phi=\langle\{S,L,D,R\},\{\pmb{b}\},S,P\rangle,
\]
where
\[
P=\{S\to\pmb{b},\ S\to LR,\ L\to LD,\ L\to D,\ D\pmb{b}\to\pmb{bb}D,\ DR\to\pmb{b}R,\ DR\to\pmb{bb}\}.
\]

Starting from $S\to LR$, repeated use of $L\to LD$ followed by $L\to D$
produces $D^jR$. The rightmost $D$ changes $R$ into one $\pmb{b}$ by
$DR\to\pmb{b}R$ (or into $\pmb{bb}$ if it is the last $D$). Each remaining $D$
then moves rightward through the current block of $\pmb{b}$'s by means of
$D\pmb{b}\to\pmb{bb}D$, doubling the length of that block. Hence each $D$
multiplies the number of $\pmb{b}$'s by $2$, and a sentential form with $j$
copies of $D$ yields exactly $\pmb{b}^{2^j}$. The first derivations are
\[
S\Rightarrow \pmb{b},
\qquad
S\Rightarrow LR\Rightarrow DR\Rightarrow \pmb{bb},
\]
and
\[
S\Rightarrow LR\Rightarrow LDR\Rightarrow DDR\Rightarrow D\pmb{b}R
\Rightarrow \pmb{bb}DR \Rightarrow \pmb{bbbb}.
\]
Therefore $L(G_\Phi)=\Phi$.
\end{enumerate}

%%------------------------------------------------------------
\item \textbf{Show $G$ (context-sensitive) is not equivalent to $G''$ in Example 8.3.}

The two grammars do \emph{not} generate the same language.

In $G''$ from Example~8.3 we have
\[
S\Rightarrow \pmb{a}SB\pmb{c}\Rightarrow \pmb{a}TB\pmb{c}
\Rightarrow \pmb{ab}T\pmb{c}\Rightarrow \pmb{abc},
\]
so
\[
\pmb{abc}\in L(G'').
\]

But $\pmb{abc}\notin L(G)$.  In the grammar $G$, after one use of
$S\to\pmb{a}SB\pmb{c}$ the only way to remove the remaining $S$ is
$S\to\lambda$, which gives
\[
\pmb{a}B\pmb{c}.
\]
That sentential form is stuck: there is no occurrence of $SB$ to which
$SB\to\pmb{b}S$ can be applied, and there is no occurrence of $\pmb{c}B$ to
which $\pmb{c}B\to B\pmb{c}$ can be applied.  Hence no derivation in $G$
produces $\pmb{abc}$.

Therefore
\[
\pmb{abc}\in L(G'')\setminus L(G),
\]
so $G$ is not equivalent to $G''$.

%%------------------------------------------------------------
\item \textbf{Context-free grammars for the given languages.}

\begin{enumerate}
\item $L_1 = \pmb{a}^*(\pmb{b}\cup\pmb{c})^* \cap \{x \mid |x|_{\pmb{a}} = |x|_{\pmb{b}} + |x|_{\pmb{c}}\}$:
Strings in $L_1$ are of the form $\pmb{a}^n\pmb{b}^j\pmb{c}^k$ where $n = j+k$:
\[
G_1: S \to \pmb{a}S\pmb{b} \mid \pmb{a}S\pmb{c} \mid \lambda
\]
Each step adds one $\pmb{a}$ at the left and one $\pmb{b}$ or $\pmb{c}$ at the right.

\item $L_2 = \{\pmb{a}^i\pmb{b}^j\pmb{c}^k \mid i+j=k\}$:
\[
G_2: S \to \pmb{a}S\pmb{c} \mid \pmb{b}T\pmb{c} \mid \lambda, \quad T \to \pmb{b}T\pmb{c} \mid \lambda
\]

\item $L_3 = \{x \in \{\pmb{a},\pmb{b},\pmb{c}\}^* \mid |x|_{\pmb{a}} + |x|_{\pmb{b}} = |x|_{\pmb{c}}\}$:
\[
G_3: S \to \pmb{a}S\pmb{c} \mid \pmb{b}S\pmb{c} \mid SS \mid \lambda
\]
\end{enumerate}

%%------------------------------------------------------------
\item \textbf{Prove $L(G') = L(G) \cup \{\lambda\}$ (Definition 8.4).}

Definition 8.4 constructs $G'$ from $G$ by adding a new start symbol $Z$ with $Z \to S \mid \lambda$ (where $S$ was the original start). Then:
$x \in L(G')$ iff $Z \DRightarrow x$ in $G'$.
$Z \to \lambda$ gives $\lambda \in L(G')$.
$Z \to S \DRightarrow x$ gives $x \in L(G) \subseteq L(G')$.
Conversely, any derivation from $Z$ begins with either $Z \to \lambda$ or $Z \to S$, so $L(G') \subseteq L(G) \cup \{\lambda\}$.
Hence $L(G') = L(G) \cup \{\lambda\}$. \qed

%%------------------------------------------------------------
\item \textbf{Prove $L(G') = L(G) \cup \{\lambda\}$ (Definition 8.6).}

Definition~8.6 constructs
\[
G'=\langle\Omega\cup\{Z\},\Sigma,Z,P\cup\{Z\to\lambda,\ Z\to S\}\rangle
\]
from the pure context-free grammar $G=\langle\Omega,\Sigma,S,P\rangle$.

If $x\in L(G)$, then $S\DRightarrow_G x$, so in $G'$ we have
\[
Z\Rightarrow S\DRightarrow x,
\]
hence $x\in L(G')$.  Also $Z\Rightarrow\lambda$, so
$\lambda\in L(G')$.  Therefore
\[
L(G)\cup\{\lambda\}\subseteq L(G').
\]

Conversely, let $x\in L(G')$.  Any derivation from the new start symbol $Z$
must begin with either $Z\to\lambda$ or $Z\to S$, because these are the only
productions having $Z$ on the left.  If the first step is $Z\to\lambda$, then
$x=\lambda$.  If the first step is $Z\to S$, the rest of the derivation uses
only productions from $P$, so $S\DRightarrow_G x$ and therefore $x\in L(G)$.
Thus
\[
L(G')\subseteq L(G)\cup\{\lambda\}.
\]

Hence
\[
L(G')=L(G)\cup\{\lambda\}.
\]
\qed

%%------------------------------------------------------------
\item \textbf{Properties of the normal form-preserving construction.}

\begin{enumerate}
\item If $G$ is right-linear (Theorem~8.8), the naive star construction $Z \to S \mid ZZ \mid \lambda$ does \emph{not} preserve right-linearity: $Z \to ZZ$ places a nonterminal in a non-final position, violating the form $A \to wB$ or $A \to w$.
A right-linear Kleene-star construction avoids $ZZ$ by folding the restart into the accepting productions of $G$: whenever $A \to w$ is a terminal production of $G$, add $A \to wS$ (to restart) and $A \to w$ (to halt). This keeps all productions right-linear.

\item Example:
\[
G_1=\langle\{S_1\},\{\pmb{a},\pmb{b}\},S_1,\{S_1\to\pmb{a}\}\rangle,\qquad
G_2=\langle\{S_2\},\{\pmb{a},\pmb{b}\},S_2,\{S_2\to\pmb{b}\}\rangle.
\]
Both grammars are in the Theorem~8.8 normal form: every production is of the
form $A\to\pmb{a}$ or $A\to\pmb{b}$, and there are no contracting productions.
But the union construction of Theorem~8.4 gives
\[
G^\cup=\langle\{Z,S_1,S_2\},\{\pmb{a},\pmb{b}\},Z,\{Z\to S_1,\ Z\to S_2,\ S_1\to\pmb{a},\ S_2\to\pmb{b}\}\rangle.
\]
The productions $Z\to S_1$ and $Z\to S_2$ are unit productions, not of the form
$A\to\pmb{a}B$ or $A\to\pmb{a}$, so $G^\cup$ is not in the normal form of
Theorem~8.8.
\item The $G^\bullet$ construction for concatenation in Example 8.18 is not normal form-preserving in general, since it requires replacing terminal productions $A \to w$ with $A \to wS_2$, which is a right-linear production but may violate other normal form constraints.
\end{enumerate}

%%------------------------------------------------------------
\item \textbf{Left-linear grammars for union, concatenation, and Kleene closure.}

Given left-linear grammars $G_1 = \langle\Omega_1,\Sigma,S_1,P_1\rangle$ and $G_2 = \langle\Omega_2,\Sigma,S_2,P_2\rangle$ with $\Omega_1 \cap \Omega_2 = \emptyset$:

\begin{enumerate}
\item \textbf{Union}: Let $H_1=G_1^r$ and $H_2=G_2^r$.  By Lemma~8.3 these are
right-linear grammars with
\[
L(H_1)=L(G_1)^r,\qquad L(H_2)=L(G_2)^r.
\]
Apply the right-linear union construction of Theorem~8.4 to obtain a
right-linear grammar $H^\cup$ for
\[
L(H_1)\cup L(H_2)=(L(G_1)\cup L(G_2))^r.
\]
Then $(H^\cup)^r$ is a left-linear grammar for $L(G_1)\cup L(G_2)$.

\item \textbf{Concatenation}: Again let $H_i=G_i^r$.  Apply the right-linear
concatenation construction to $H_2$ and $H_1$ (in that order) to obtain a
right-linear grammar $H^\bullet$ such that
\[
L(H^\bullet)=L(H_2)\cdot L(H_1)=L(G_2)^rL(G_1)^r=(L(G_1)\cdot L(G_2))^r.
\]
Then $(H^\bullet)^r$ is a left-linear grammar for $L(G_1)\cdot L(G_2)$.

\item \textbf{Kleene closure}: Let $H_1=G_1^r$, and apply the right-linear
Kleene-star construction of Theorem~8.5 to get a right-linear grammar $H_*$ for
\[
L(H_*)=L(H_1)^*=(L(G_1)^*)^r.
\]
Then $(H_*)^r$ is a left-linear grammar for $L(G_1)^*$.
\end{enumerate}

\end{enumerate}
